% Length of (1, 2, -2) from the Pythagorean theorem twice: first on the floor, then with the height.
\begin{tikzpicture}[scale=1.45, x={(-0.55cm,-0.42cm)}, y={(1cm,0cm)}, z={(0cm,1cm)}]
\draw[siv, ->] (0,0,0) -- (3.6,0,0) node[below left]{$x$};
\draw[siv, ->] (0,0,0) -- (0,5.6,0) node[right]{$y$};
\draw[siv, ->] (0,0,0) -- (0,0,2.8) node[above]{$z$};
\fill[siv, opacity=0.15] (2,3,0) -- (3,3,0) -- (3,5,0) -- cycle;
\fill[acc, opacity=0.10] (2,3,2) -- (2,3,0) -- (3,5,0) -- cycle;
\draw (2,3,0) -- node[above left]{$1$} (3,3,0) -- node[below]{$2$} (3,5,0);
\draw (2,3,2) -- node[left, pos=0.25]{$2$} (2,3,0);
\draw[siv, line width=0.8pt] (2,3,0) -- node[above, pos=0.6]{$\sqrt 5$} (3,5,0);
\draw[->, acc, line width=1pt] (2,3,2) -- node[above right, pos=0.3]{$3$} (3,5,0);
\draw[siv, line width=0.45pt] (2.7,3,0) -- (2.7,3.3,0) -- (3,3.3,0);
\draw[siv, line width=0.45pt] (2,3,0.3) -- (2.134,3.268,0.3) -- (2.134,3.268,0);
\node[dot] at (2,3,2) {};
\node[above] at (2,3,2) {$P$};
\node[right] at (3,5,0) {$Q$};
\end{tikzpicture}
